\textbf{Hardware Platforms.}
We evaluate \NAME{} on an NVIDIA GH200 Grace Hopper Superchip (H100 GPU, Hopper)
and an NVIDIA RTX 5090 GPU (Blackwell).

\textbf{Workloads.}
We evaluate the attention operators of five transformer models
(Appendix Table~\ref{tab:operators}):
global attention from ViT 1.2B, causal attention (MHA) from Llama 2 7B, Qwen2 7B, and Qwen3 8B,
and grouped-query attention (GQA) from GLM-4 9B.
In end-to-end experiments, we run each model in prefill mode
on batch sizes 1 and 8, sequence lengths $2^{10}, \dots, 2^{15}$,
and FP16 and FP8-E4M3 precisions;
in operator experiments, we evaluate the three attention operators individually
on the same inputs.
We refer to the tuple of hardware platform, operator, input shape, and input precision as
a \textbf{setup}.
For each setup, we run \NAME{}'s auto-tuner for 256 empirical measurements.

\textbf{Baselines.}
We compare \NAME{} to 12 baselines (names we use are underlined):
seven compilers, namely \uline{Tawa}~\cite{tawa} and \uline{TileLang}~\cite{tilelang}
with their FlashAttention-3-style attention kernels,
two Triton implementations (\uline{OpenAI Triton}~\cite{openai-fused-attn}
and \uline{Tri-Dao Triton}~\cite{flash-attn-repo}), \uline{Helion},
FlexAttention (\uline{FlexAttn}) in PyTorch, and \uline{TVM};
and five hand-optimized libraries, namely \uline{cuDNN},
Dao-AILab's CUTLASS FlashAttention (\uline{Tri-Dao})~\cite{flash-attn-repo},
\uline{FlashInfer}~\cite{flashinfer}, \uline{ThunderKittens}~\cite{thunderkittens},
and \uline{PyTorch}'s built-in attention~\cite{pytorch}.
Appendix Table~\ref{tab:baselines} lists versions and configurations.
We do not evaluate FlashAttention-4~\cite{zadouri2026flashattention4algorithmkernelpipelining}
because it uses approximation techniques outside the scope of exact attention.

\textbf{Profiling and Performance Reporting.}
We report throughput in TFLOPs per second, dividing the theoretical attention FLOPs
of each configuration~\cite{Megatron-LM} by the measured latency,
which we measure with CUDA Events in end-to-end experiments
and with Nsight Systems in operator experiments (Appendix~\ref{sec:exp-details}).
We warm up once, report the mean latency of 10 measurements,
and average across setups with geometric means.

% \yuchen{GH200 has H100 as GPU} \yifan{thanks for the correction!}
%\sasa{Motivate that these are sota hw architectures with modern features...}
% \NA{Sparse} & \NA{Deepseek Sparse Attention} & \NA{DeepSeek-V3.2} \\
%\sasa{Can we motivate why these networks are representative?}
% Each operator is evaluated in both prefill and decoding modes:
% in prefill (PF), the query sequence length $s_q$ equals the key/value sequence length $s_{kv}$,
% and in decoding (DC), $s_q = 1$.
% Variations on masking (Global and Causal) collapse into the same operator under decoding
% mode, so we select one.
% Each operator is evaluated in prefill (PF) mode, where the query sequence length $s_q$
% equals the key/value sequence length $s_{kv}$.
%\yifan{Reminder: ThunderKittens kernel is not ``official'' but provided by one of the maintainers.}
% \sasa{Doublechk; cite the method if we have from somewhere}
%We also report \textbf{relative performance} of a baseline relative to \NAME{},
%which is $t_0 / t_i$ for the $i$-th baseline.
